% Supplementary Note S-Tooling: Comparative Positioning in the Analog CIM Ecosystem
% Date: 2026-04-25

\section*{Supplementary Note S-Tooling: Tooling Positioning in the Analog CIM Ecosystem}
\label{supp:tooling}

This note situates the present framework within the broader analog compute-in-memory simulation and training ecosystem. The goal is honest comparative positioning: our work is complementary to (not replacing) existing mature tools, and we cite conceptual lineage where due.

\subsection*{S-T.1 Quantitative cross-check}

We performed a controlled cross-comparison on a 1000-image deterministic CIFAR-10 subset:

\begin{itemize}
    \item \textbf{Clean baseline:} Our framework reports 86.2\%; CrossSim reports 83.7\% at 8-bit ADC.
    \item \textbf{Noisy 3-run Monte Carlo} ($\sigma_{\text{D2D}} = 5\%$): We report 81.63$\pm$0.56\%; CrossSim reports 67.20$\pm$2.67\%.
\end{itemize}

The 14.4~pp divergence reflects a physical-modeling convention difference: we use multiplicative D2D mismatch (relative to conductance range), whereas CrossSim's default uses additive mismatch (absolute conductance deviation). Both conventions are internally consistent; the divergence is a mapping-convention artifact, not a correctness discrepancy. CrossSim is substantially slower per-evaluation for full parameter sweeps, trading physical fidelity for sweep speed.

\subsection*{S-T.2 Conceptual ancestry}

We acknowledge AIHWKit (Rasch \emph{et al.}\ 2023) as the conceptual ancestor of the train-with-differentiable-surrogate / eval-with-ADC discipline and the per-epoch noise resampling primitives used in the present framework.

\subsection*{S-T.3 AIHWKit head-to-head attempt}

We attempted a direct head-to-head replication of the canonical Tiny-ViT V4 checkpoint in AIHWKit~v0.9.1. Conversion failed with \texttt{ValueError: Only one group is supported}, because Tiny-ViT's patch-embedding and stage-0 blocks employ grouped convolutions (depthwise-separable filters) that AIHWKit's \texttt{AnalogConv2d} does not currently map. A ResNet-18 equivalent (no grouped convolutions) converts successfully, but that architecture is outside the present study's scope. This limitation underscores why a custom framework was necessary for Vision-Transformer CIM evaluation: contemporary open-source analog-training toolkits are optimized for standard convolutions and fully-connected layers, not for the heterogeneous operator mix (grouped conv, attention, MLP) found in modern ViTs.
